% ==========================================================================
% HELIX Platform -- PoC Engineering Roadmap
% LaTeX Document
% ==========================================================================
\documentclass[11pt,a4paper]{report}

\usepackage[T1]{fontenc}
\usepackage[utf8]{inputenc}
\usepackage[english]{babel}
\usepackage{geometry}
\usepackage{setspace}
\usepackage{microtype}
\usepackage{parskip}
\usepackage{xcolor}
\usepackage{booktabs}
\usepackage{array}
\usepackage{float}
\usepackage{fancyhdr}
\usepackage{titlesec}
\usepackage{enumitem}
\usepackage{tcolorbox}
\usepackage{hyperref}
\usepackage{longtable}
\usepackage{graphicx}
\usepackage{listings}
\usepackage{etoolbox}
\usepackage{tabularx}
\usepackage{multirow}
\usepackage{makecell}

\geometry{a4paper, top=2cm, bottom=2cm, left=2cm, right=2cm, headheight=15pt}

% ── Colors ───────────────────────────────────────────────────────────────
\definecolor{helixpurple}{HTML}{3C3489}
\definecolor{helixpurplelight}{HTML}{EEEDFE}
\definecolor{helixpurplemid}{HTML}{5A50B0}
\definecolor{helixteal}{HTML}{0F6E56}
\definecolor{helixteallight}{HTML}{E1F5EE}
\definecolor{helixamber}{HTML}{7A4A00}
\definecolor{helixamberlight}{HTML}{FDF3DC}
\definecolor{helixgray}{HTML}{3A3A38}
\definecolor{helixgraylight}{HTML}{F4F3EF}
\definecolor{helixborder}{HTML}{CCCAC0}
\definecolor{textsecondary}{HTML}{5A5958}
\definecolor{codebg}{HTML}{F8F7F4}
\definecolor{codeframe}{HTML}{CCCAC0}

% ── tcolorbox ────────────────────────────────────────────────────────────
\tcbuselibrary{skins,breakable,listings}

\newtcolorbox{phasebox}[2][]{
    enhanced, breakable,
    colback=helixpurplelight, colframe=helixpurple,
    fonttitle=\bfseries\small\color{white},
    coltitle=white,
    attach boxed title to top left={yshift=-2mm, xshift=6pt},
    boxed title style={colback=helixpurple, colframe=helixpurple, arc=3pt},
    title=#2,
    left=8pt, right=8pt, top=10pt, bottom=8pt,
    arc=5pt, boxrule=0.6pt, #1
}

\newtcolorbox{taskbox}[1][]{
    enhanced, breakable,
    colback=white, colframe=helixborder,
    left=8pt, right=8pt, top=6pt, bottom=6pt,
    arc=4pt, boxrule=0.4pt, #1
}

\newtcolorbox{infobox}[1][]{
    enhanced, breakable,
    colback=helixteallight, colframe=helixteal,
    fonttitle=\bfseries\small\color{helixteal},
    left=8pt, right=8pt, top=6pt, bottom=6pt,
    arc=4pt, boxrule=0.6pt, #1
}

\newtcolorbox{warnbox}[1][]{
    enhanced, breakable,
    colback=helixamberlight, colframe=helixamber,
    fonttitle=\bfseries\small\color{helixamber},
    left=8pt, right=8pt, top=6pt, bottom=6pt,
    arc=4pt, boxrule=0.6pt, #1
}

% ── Listings (code blocks) ───────────────────────────────────────────────
\lstdefinestyle{helixcode}{
    backgroundcolor=\color{codebg},
    frame=single,
    rulecolor=\color{codeframe},
    framerule=0.4pt,
    basicstyle=\ttfamily\small,
    keywordstyle=\color{helixpurplemid}\bfseries,
    commentstyle=\color{textsecondary}\itshape,
    stringstyle=\color{helixteal},
    numbers=left,
    numberstyle=\tiny\color{helixborder},
    numbersep=5pt,
    breaklines=true,
    breakatwhitespace=false,
    showstringspaces=false,
    tabsize=4,
    xleftmargin=12pt,
    xrightmargin=4pt,
}

\lstset{style=helixcode}

\lstdefinelanguage{JavaScript}{
    keywords={typeof, new, true, false, catch, function, return, null, catch, switch, var, if, in, while, do, else, case, break, const, let, async, await, export, default, import, from, class, extends, super, this, yield},
    keywordstyle=\color{helixpurplemid}\bfseries,
    ndkeywords={=, =>, ===, !=, <, >, <=, >=, +=, -=, *=, /=},
    sensitive=true,
    comment=[l]{//},
    morecomment=[s]{/*}{*/},
    morestring=[b]',
    morestring=[b]",
    morestring=[b]`,
}

\lstdefinelanguage{JSON}{
    keywords={true, false, null},
    keywordstyle=\color{helixpurplemid}\bfseries,
    sensitive=true,
    morestring=[b]",
    morecomment=[l]{//},
}

% ── Section styling ──────────────────────────────────────────────────────
\titleformat{\chapter}[block]
{\normalfont\LARGE\bfseries\color{helixpurple}}
{\thechapter.}{12pt}{}
[\vspace{3pt}\color{helixborder}\rule{\linewidth}{0.4pt}\vspace{6pt}]

\titleformat{\section}
{\normalfont\large\bfseries\color{helixgray}}
{\thesection}{10pt}{}
[\vspace{2pt}\color{helixborder}\rule{\linewidth}{0.3pt}]

\titleformat{\subsection}
{\normalfont\normalsize\bfseries\color{helixpurplemid}}
{\thesubsection}{8pt}{}

\titlespacing*{\chapter}{0pt}{10pt}{22pt}
\titlespacing*{\section}{0pt}{16pt}{8pt}
\titlespacing*{\subsection}{0pt}{12pt}{6pt}

% ── Header / Footer ──────────────────────────────────────────────────────
\pagestyle{fancy}
\fancyhf{}
\renewcommand{\headrulewidth}{0.4pt}
\renewcommand{\footrulewidth}{0.4pt}
\fancyhead[L]{\small\color{textsecondary}\scshape{HELIX Platform}}
\fancyhead[R]{\small\color{textsecondary}PoC Engineering Roadmap}
\fancyfoot[C]{\small\color{textsecondary}\thepage}
\fancyfoot[R]{\small\color{textsecondary}2025--2026}

% ── Custom column types ──────────────────────────────────────────────────
\newcolumntype{L}[1]{>{\raggedright\arraybackslash}p{#1}}
\newcolumntype{C}[1]{>{\centering\arraybackslash}p{#1}}
\newcolumntype{R}[1]{>{\raggedleft\arraybackslash}p{#1}}

% ── List styling ─────────────────────────────────────────────────────────
\setlist[itemize]{leftmargin=1.5em, itemsep=2pt, topsep=3pt}
\setlist[enumerate]{leftmargin=1.7em, itemsep=2pt, topsep=3pt}

% ── Hyperref ─────────────────────────────────────────────────────────────
\hypersetup{
    colorlinks=true,
    linkcolor=helixpurple,
    urlcolor=helixteal,
    citecolor=helixgray,
    pdftitle={HELIX -- PoC Engineering Roadmap},
    pdfauthor={Yasseene},
    pdfsubject={Cybersecurity Operations Platform},
}

% ── Custom commands ──────────────────────────────────────────────────────
\newcommand{\todo}{\textcolor{helixamber}{\textbf{[ ]}}}
\newcommand{\done}{\textcolor{helixteal}{\textbf{[x]}}}
\newcommand{\inprogress}{\textcolor{helixpurplemid}{\textbf{[~]}}}
\newcommand{\blocked}{\textcolor{red}{\textbf{[!]}}}

\newcommand{\filepath}[1]{\texttt{\small#1}}

% ==========================================================================
\begin{document}
\setstretch{1.15}

% ── Cover ──────────────────────────────────────────────────────────────────
\begin{titlepage}
    \thispagestyle{empty}
    \noindent\colorbox{helixpurple}{%
        \parbox{\dimexpr\textwidth+2\fboxsep\relax}{%
            \vspace{8pt}
            \hspace{8pt}\textcolor{white}{\small\scshape{Proof of Concept Engineering Roadmap}}
            \vspace{8pt}
        }%
    }
    \vspace{3cm}
    \begin{center}
        {\fontsize{52}{60}\selectfont\bfseries\color{helixpurple} HELIX}\\[10pt]
        {\Large\color{textsecondary}Cybersecurity Operations Platform}\\[8pt]
        \color{helixborder}\rule{0.65\textwidth}{0.5pt}\\[30pt]
        {\large\color{textsecondary}PoC Engineering Roadmap}\\[6pt]
        {\normalsize\color{textsecondary}Learn \& Build Hybrid Approach}
    \end{center}
    \vfill
    \begin{center}
        \begin{tabular}{L{5.2cm} L{7.5cm}}
            \toprule
            \textbf{Document Type}    & Engineering Roadmap \\
            \textbf{Project Phase}    & Proof of Concept (PoC) \\
            \textbf{Document Version} & 2.0 \\
            \textbf{Academic Year}    & 2025 -- 2026 \\
            \textbf{Author}           & Yasseene \\
            \textbf{Classification}   & Academic -- Confidential \\
            \bottomrule
        \end{tabular}
    \end{center}
    \vspace{1.5cm}
    \noindent\colorbox{helixgraylight}{%
        \parbox{\dimexpr\textwidth+2\fboxsep\relax}{%
            \vspace{8pt}
            \hspace{8pt}\small\color{textsecondary}
            Tek-Up University \hfill Ariana, Tunisia \hfill Academic Year 2025--2026
            \vspace{8pt}
        }%
    }
\end{titlepage}

{
    \hypersetup{linkcolor=black}
    \tableofcontents
    \thispagestyle{empty}
}
\clearpage

% ══════════════════════════════════════════════════════════════════════════
% INTRODUCTION
% ══════════════════════════════════════════════════════════════════════════
\chapter*{Introduction}
\addcontentsline{toc}{chapter}{Introduction}

This document is your personal engineering execution blueprint, adapted to your actual project architecture: a decoupled SPA frontend, a layered PHP API (Controller $\rightarrow$ Service $\rightarrow$ Model), a Python ML module, and a Java utility module. Every file reference in this roadmap maps to a real file in your tree. Follow it sequentially, create files only when you reach them, and mark your progress as you go.

\section*{How to Use This Roadmap}

\begin{center}
\begin{tabular}{C{2cm} L{12cm}}
    \toprule
    \done\ Not started & Mark as completed when done \\
    \inprogress\ In progress & Currently working on \\
    \todo\ Completed & Task finished \\
    \blocked\ Blocked & Needs review or unblocking \\
    \bottomrule
\end{tabular}
\end{center}

\begin{infobox}[title=Core Principle]
You do not move to the next phase until the current phase's \textbf{Definition of Done} is satisfied. You create files only when you need them -- not before.
\end{infobox}

% ══════════════════════════════════════════════════════════════════════════
% ARCHITECTURE REFERENCE
% ══════════════════════════════════════════════════════════════════════════
\chapter*{Architecture Reference}
\addcontentsline{toc}{chapter}{Architecture Reference}

Before starting, internalize this flow. Every feature you build will follow it:

\begin{verbatim}
Browser (frontend/)
    |
    |  HTTP fetch() via api.js
    v
api/index.php          <- Front Controller: parses route, applies middleware, dispatches
    |
    |-- middleware/AuthMiddleware.php     <- Is the user logged in?
    |-- middleware/RoleMiddleware.php     <- Is the user allowed?
    |
    v
api/controllers/XController.php         <- Parse request, call service, return response
    |
    v
api/services/XService.php               <- All business logic lives here
    |
    v
api/models/X.php                        <- All DB queries live here (PDO)
    |
    v
MySQL (helix_db)
\end{verbatim}

\textbf{Key architectural decisions already made:}

\begin{itemize}
    \item The frontend is completely decoupled -- it never touches PHP directly
    \item \filepath{api/index.php} is the single entry point -- all routes go through it
    \item \filepath{api/config/Cors.php} handles CORS because frontend and API are on different paths
    \item \filepath{api/helpers/Response.php} standardizes all JSON responses
    \item \filepath{api/helpers/Validator.php} centralizes input validation
    \item \filepath{api/services/JavaBridge.php} is the only place that touches \texttt{shell\_exec()}
\end{itemize}

% ══════════════════════════════════════════════════════════════════════════
% GLOBAL PHASE MAP
% ══════════════════════════════════════════════════════════════════════════
\chapter*{Global Phase Map}
\addcontentsline{toc}{chapter}{Global Phase Map}

\begin{verbatim}
PHASE 0 -> PHASE 1 -> PHASE 2 -> PHASE 3 -> PHASE 4 -> PHASE 5 -> PHASE 6 -> PHASE 7 -> PHASE 8 -> PHASE 9 -> PHASE 10
  |           |           |           |           |           |           |           |           |          |           |
Setup     Architecture  Auth +    Dashboard   Log         Alert       ML          AI        Scanner  Integration  Docker
& Git     & DB Design   RBAC      & Routing  Ingestion   Module     Detection  Assistant    Module   & Testing   & Deploy
\end{verbatim}

\textbf{Dependency chain (strict):}

\begin{itemize}
    \item Phase 1 requires Phase 0
    \item Phase 2 requires Phase 1 (DB schema + routing foundation in place)
    \item Phase 3 requires Phase 2 (auth middleware working, frontend routing exists)
    \item Phase 4 requires Phase 3 (dashboard shell exists as UI home)
    \item Phase 5 requires Phase 4 (log data is the raw input for alerts + ML)
    \item Phase 6 requires Phase 5 (alerts are generated from scored logs)
    \item Phase 7 requires Phase 5 (ML needs log data to score)
    \item Phase 8 requires Phase 6 (AI analyzes alerts enriched by ML scores)
    \item Phase 9 requires Phase 2 (scanner only needs auth, independent of logs/ML)
    \item Phase 10 requires Phases 2--9 (all modules complete and individually tested)
    \item Phase 11 requires Phase 10 (containerize only a stable application)
\end{itemize}

% ══════════════════════════════════════════════════════════════════════════
% PHASE 0
% ══════════════════════════════════════════════════════════════════════════
\chapter{PHASE 0 -- Project Setup, Git Workflow \& Development Environment}

\begin{phasebox}{Phase 0: Setup \& Environment}
\textbf{Why this comes first:} Without a reproducible environment and version control, everything you build is fragile. One wrong XAMPP config, one missing \filepath{.gitignore}, and you lose hours. This phase costs 2 hours and saves 20.
\end{phasebox}

\section*{Objective}
Establish the complete development environment, Git discipline, Apache serving strategy, Composer setup, and Python environment -- before writing a single line of application logic.

\section*{Learning Goals}
\begin{itemize}
    \item Understand Git branching strategy for solo development
    \item Understand how Apache serves two directories under one vhost (frontend + api)
    \item Understand Composer as PHP's dependency manager and autoloader
    \item Understand \filepath{pyproject.toml} vs \filepath{requirements.txt} (and why you're using the former)
    \item Understand the Front Controller pattern (\filepath{api/index.php})
\end{itemize}

\section*{Key Theoretical Concepts}
\begin{itemize}
    \item \textbf{Git Flow (simplified):} \texttt{main} = stable releases, \texttt{dev} = integration branch, \texttt{feature/*} = active work
    \item \textbf{Apache vhost + Alias directive:} How to serve \filepath{frontend/} at \texttt{/} and \filepath{api/} at \texttt{/api} from one domain
    \item \textbf{Composer autoloading (PSR-4):} How \texttt{use App\textbackslash Controllers\textbackslash AuthController} resolves to a file path automatically
    \item \textbf{Front Controller:} One entry point that routes all requests -- no 30 separate PHP files each handling their own routing
    \item \textbf{CORS (Cross-Origin Resource Sharing):} Why the browser blocks your \texttt{fetch('/api/...')} without the right headers
\end{itemize}

\section{0.1 -- Git Repository Setup}

\begin{itemize}
    \item[\done] Create private repository on GitHub/GitLab
    \item[\done] Clone locally into your project root
    \item[\done] Create \filepath{.gitignore}:
\end{itemize}

\begin{lstlisting}[language=bash]
# PHP
vendor/
composer.lock (optional)

# Python
python-module/.venv/
python-module/__pycache__/
python-module/models/*.joblib

# Java
java-module/build/*.class

# Environment
api/config/.env
.env

# IDE
.vscode/
.idea/
*.DS_Store
\end{lstlisting}

\begin{itemize}
    \item[\todo] Create \filepath{CHANGELOG.md} -- update at end of every phase
    \item[\todo] Initial commit: \texttt{git commit -m "chore: initialize HELIX project"}
    \item[\todo] Create \texttt{dev} branch: \texttt{git checkout -b dev}
\end{itemize}

\textbf{Branching strategy:}
\begin{itemize}
    \item Work on \texttt{feature/*} branches cut from \texttt{dev}
    \item Merge \texttt{feature/*} $\rightarrow$ \texttt{dev} when a phase is complete and tested
    \item Merge \texttt{dev} $\rightarrow$ \texttt{main} only when a full milestone works end-to-end
\end{itemize}

\section{0.2 -- Apache Vhost Configuration (Fedora/RHEL)}

This is the critical decision from the review. You have \filepath{frontend/} and \filepath{api/} as siblings. Define how Apache serves both under \texttt{helix.local}. On Fedora, Apache config lives in \filepath{/etc/httpd/}.

\subsection*{Step 1: Install Apache, PHP, and dependencies}

\begin{lstlisting}[language=bash]
sudo dnf install httpd php php-cli
sudo dnf install php-pdo php-pdo_mysql php-json
sudo systemctl start httpd
sudo systemctl enable httpd
\end{lstlisting}

Verify Apache is running: \texttt{curl http://localhost}

\subsection*{Step 2: Create vhost configuration}

Create \filepath{/etc/httpd/conf.d/helix.conf} (requires \texttt{sudo}):

\begin{lstlisting}[language=bash]
<VirtualHost *:80>
    ServerName helix.local
    DocumentRoot /home/yasseene/HELIX/frontend

    <Directory /home/yasseene/HELIX/frontend>
        Options Indexes FollowSymLinks
        AllowOverride All
        Require all granted

        <IfModule mod_rewrite.c>
            RewriteEngine On
            RewriteBase /
            RewriteCond %{REQUEST_FILENAME} !-f
            RewriteCond %{REQUEST_FILENAME} !-d
            RewriteRule ^(.*)$ index.html [QSA,L]
        </IfModule>
    </Directory>

    Alias /api /home/yasseene/HELIX/api

    <Directory /home/yasseene/HELIX/api>
        Options FollowSymLinks
        AllowOverride All
        Require all granted

        <IfModule mod_rewrite.c>
            RewriteEngine On
            RewriteBase /api
            RewriteCond %{REQUEST_FILENAME} !-f
            RewriteCond %{REQUEST_FILENAME} !-d
            RewriteRule ^(.*)$ index.php?path=$1 [QSA,L]
        </IfModule>
    </Directory>

    <FilesMatch \.php$>
        SetHandler application/x-httpd-php
    </FilesMatch>

    ErrorLog /var/log/httpd/helix_error.log
    CustomLog /var/log/httpd/helix_access.log combined
</VirtualHost>
\end{lstlisting}

\subsection*{Step 3: Enable mod\_rewrite}

On Fedora/RHEL, \texttt{mod\_rewrite} is usually enabled by default. Verify:

\begin{lstlisting}[language=bash]
grep -i "rewrite" /etc/httpd/conf.modules.d/*.conf
\end{lstlisting}

\subsection*{Step 4: Add to system hosts file}

\begin{lstlisting}[language=bash]
echo "127.0.0.1    helix.local" | sudo tee -a /etc/hosts
ping helix.local
\end{lstlisting}

\subsection*{Step 5: Fix file permissions}

\begin{lstlisting}[language=bash]
sudo usermod -a -G $(id -gn) apache
chmod 755 /home/yasseene/HELIX
chmod 755 /home/yasseene/HELIX/{frontend,api}
\end{lstlisting}

\subsection*{Step 6: Validate and restart Apache}

\begin{lstlisting}[language=bash]
sudo httpd -t
sudo systemctl restart httpd
sudo systemctl status httpd
\end{lstlisting}

\textbf{Consequence for \filepath{api.js}:} All frontend API calls will use:

\begin{lstlisting}[language=JavaScript]
const API_BASE = "/api"; // relative -- works on same origin
\end{lstlisting}

Never hardcode \texttt{http://localhost/...} -- relative paths survive domain changes.

\begin{itemize}
    \item[\done] Install Apache, PHP, and dependencies
    \item[\done] Create \filepath{/etc/httpd/conf.d/helix.conf} with vhost above
    \item[\done] Add \texttt{127.0.0.1 helix.local} to \filepath{/etc/hosts}
    \item[\done] Fix permissions
    \item[\done] Verify mod\_rewrite
    \item[\done] Run \texttt{sudo httpd -t} (syntax check)
    \item[\done] Run \texttt{sudo systemctl restart httpd}
    \item[\done] Verify: \texttt{http://helix.local} serves \filepath{frontend/index.html}
    \item[\done] Create \filepath{api/index.php} with \texttt{echo "API OK";} temporarily
    \item[\done] Verify: \texttt{http://helix.local/api} returns "API OK"
    \item[\done] Check logs: \texttt{tail -f /var/log/httpd/helix\_error.log} if anything fails
\end{itemize}

\section{0.3 -- MySQL Setup}

\begin{itemize}
    \item[\done] Start MySQL in XAMPP
    \item[\done] Create database: \texttt{helix\_db}
    \item[\done] Create dedicated user with limited privileges:
\end{itemize}

\begin{lstlisting}[language=SQL]
CREATE USER 'helix_user'@'localhost' IDENTIFIED BY 'your_password';
GRANT SELECT, INSERT, UPDATE, DELETE ON helix_db.* TO 'helix_user'@'localhost';
FLUSH PRIVILEGES;
\end{lstlisting}

\begin{itemize}
    \item[\todo] Verify connection via phpMyAdmin
\end{itemize}

\section{0.4 -- PHP Composer Setup}

\filepath{composer.json}:

\begin{lstlisting}[language=JSON]
{
  "name": "helix/api",
  "description": "HELIX Cybersecurity Platform API",
  "type": "project",
  "require": {
    "php": ">=8.1"
  },
  "require-dev": {
    "phpunit/phpunit": "^10.0"
  },
  "autoload": {
    "psr-4": {
      "App\\": "api/"
    }
  },
  "autoload-dev": {
    "psr-4": {
      "Tests\\": "tests/php/"
    }
  }
}
\end{lstlisting}

\begin{itemize}
    \item[\todo] Run: \texttt{composer install}
    \item[\todo] Verify \filepath{vendor/autoload.php} was created
    \item[\todo] Add \texttt{require\_once \_\_DIR\_\_ . '/../../vendor/autoload.php';} to \filepath{api/index.php}
\end{itemize}

\textbf{Why PSR-4 matters:} With this config, \texttt{use App\textbackslash Controllers\textbackslash AuthController} automatically maps to \filepath{api/controllers/AuthController.php}. You never write manual \texttt{require} statements for your own classes.

\section{0.5 -- Environment Configuration}

\filepath{api/config/Environment.php}:

\begin{lstlisting}[language=PHP]
<?php
namespace App\Config;

class Environment
{
    private static array $vars = [];

    public static function load(string $envFile): void
    {
        if (!file_exists($envFile)) return;

        foreach (file($envFile, FILE_IGNORE_NEW_LINES | FILE_SKIP_EMPTY_LINES) as $line) {
            if (str_starts_with(trim($line), '#')) continue;
            [$key, $value] = explode('=', $line, 2);
            self::$vars[trim($key)] = trim($value);
            $_ENV[trim($key)] = trim($value);
        }
    }

    public static function get(string $key, mixed $default = null): mixed
    {
        return self::$vars[$key] ?? $_ENV[$key] ?? getenv($key) ?: $default;
    }
}
\end{lstlisting}

Create \filepath{api/config/.env} (gitignored):

\begin{lstlisting}[language=bash]
DB_HOST=localhost
DB_NAME=helix_db
DB_USER=helix_user
DB_PASS=your_password
OPENAI_API_KEY=sk-...
SESSION_LIFETIME=3600
PYTHON_PATH=C:/path/to/.venv/Scripts/python.exe
PYTHON_SCRIPT=python-module/anomaly_detector.py
JAVA_BUILD_PATH=java-module/build
\end{lstlisting}

Create \filepath{api/config/.env.example} (committed to git -- no real values):

\begin{lstlisting}[language=bash]
DB_HOST=localhost
DB_NAME=helix_db
DB_USER=helix_user
DB_PASS=
OPENAI_API_KEY=
SESSION_LIFETIME=3600
PYTHON_PATH=
PYTHON_SCRIPT=python-module/anomaly_detector.py
JAVA_BUILD_PATH=java-module/build
\end{lstlisting}

\section{0.6 -- Database Config}

\filepath{api/config/Database.php}:

\begin{lstlisting}[language=PHP]
<?php
namespace App\Config;

use PDO;
use PDOException;

class Database
{
    private static ?PDO $instance = null;

    public static function getInstance(): PDO
    {
        if (self::$instance === null) {
            $host = Environment::get('DB_HOST', 'localhost');
            $name = Environment::get('DB_NAME');
            $user = Environment::get('DB_USER');
            $pass = Environment::get('DB_PASS');

            $dsn = "mysql:host={$host};dbname={$name};charset=utf8mb4";
            $options = [
                PDO::ATTR_ERRMODE            => PDO::ERRMODE_EXCEPTION,
                PDO::ATTR_DEFAULT_FETCH_MODE => PDO::FETCH_ASSOC,
                PDO::ATTR_EMULATE_PREPARES   => false,
            ];

            try {
                self::$instance = new PDO($dsn, $user, $pass, $options);
            } catch (PDOException $e) {
                error_log("DB connection failed: " . $e->getMessage());
                http_response_code(500);
                echo json_encode(['error' => 'Database unavailable']);
                exit;
            }
        }

        return self::$instance;
    }
}
\end{lstlisting}

\begin{itemize}
    \item[\todo] Test: add a \texttt{Database::getInstance();} call to \filepath{api/index.php} temporarily $\rightarrow$ no error
\end{itemize}

\section{0.7 -- CORS Configuration}

\filepath{api/config/Cors.php}:

\begin{lstlisting}[language=PHP]
<?php
namespace App\Config;

class Cors
{
    public static function handle(): void
    {
        $allowedOrigins = ['http://helix.local', 'http://localhost'];
        $origin = $_SERVER['HTTP_ORIGIN'] ?? '';

        if (in_array($origin, $allowedOrigins, true)) {
            header("Access-Control-Allow-Origin: {$origin}");
        }

        header('Access-Control-Allow-Methods: GET, POST, PUT, DELETE, OPTIONS');
        header('Access-Control-Allow-Headers: Content-Type, Authorization, X-Requested-With');
        header('Access-Control-Allow-Credentials: true');

        if ($_SERVER['REQUEST_METHOD'] === 'OPTIONS') {
            http_response_code(204);
            exit;
        }
    }
}
\end{lstlisting}

\section{0.8 -- Helper Classes}

\filepath{api/helpers/Response.php}:

\begin{lstlisting}[language=PHP]
<?php
namespace App\Helpers;

class Response
{
    public static function json(mixed $data, int $status = 200): void
    {
        http_response_code($status);
        header('Content-Type: application/json');
        echo json_encode($data);
        exit;
    }

    public static function success(mixed $data = null, string $message = 'OK', int $status = 200): void
    {
        self::json(['success' => true, 'message' => $message, 'data' => $data], $status);
    }

    public static function error(string $message, int $status = 400): void
    {
        self::json(['success' => false, 'error' => $message], $status);
    }

    public static function unauthorized(): void
    {
        self::error('Unauthorized', 401);
    }

    public static function forbidden(): void
    {
        self::error('Forbidden', 403);
    }
}
\end{lstlisting}

\filepath{api/helpers/Validator.php}:

\begin{lstlisting}[language=PHP]
<?php
namespace App\Helpers;

class Validator
{
    private array $errors = [];

    public function required(string $field, mixed $value): self
    {
        if (empty($value) && $value !== '0') {
            $this->errors[$field] = "{$field} is required";
        }
        return $this;
    }

    public function email(string $field, mixed $value): self
    {
        if (!filter_var($value, FILTER_VALIDATE_EMAIL)) {
            $this->errors[$field] = "{$field} must be a valid email";
        }
        return $this;
    }

    public function minLength(string $field, mixed $value, int $min): self
    {
        if (strlen((string)$value) < $min) {
            $this->errors[$field] = "{$field} must be at least {$min} characters";
        }
        return $this;
    }

    public function passes(): bool
    {
        return empty($this->errors);
    }

    public function errors(): array
    {
        return $this->errors;
    }
}
\end{lstlisting}

\section{0.9 -- Front Controller Bootstrap}

\filepath{api/index.php} -- the single entry point for all API calls:

\begin{lstlisting}[language=PHP]
<?php
declare(strict_types=1);

require_once __DIR__ . '/../vendor/autoload.php';

use App\Config\Environment;
use App\Config\Cors;
use App\Helpers\Response;

// 1. Load environment
Environment::load(__DIR__ . '/config/.env');

// 2. Handle CORS (must be before any output)
Cors::handle();

// 3. Parse route
$method = $_SERVER['REQUEST_METHOD'];
$uri    = parse_url($_SERVER['REQUEST_URI'], PHP_URL_PATH);

// Strip /api prefix (added by Apache Alias)
$uri = preg_replace('#^/api#', '', $uri);
$uri = rtrim($uri, '/') ?: '/';

// 4. Route dispatch (will grow with each phase)
$routes = [];

// Each phase will add routes here:
// $routes['POST']['/auth/login'] = [AuthController::class, 'login'];

// 5. Dispatch
if (isset($routes[$method][$uri])) {
    [$controllerClass, $action] = $routes[$method][$uri];
    $controller = new $controllerClass();
    $controller->$action();
} else {
    Response::error('Route not found', 404);
}
\end{lstlisting}

\begin{itemize}
    \item[\todo] Verify: \texttt{http://helix.local/api/nonexistent} returns \texttt{\{"success":false,"error":"Route not found"\}}
\end{itemize}

\section{0.10 -- Python Environment}

\begin{itemize}
    \item[\todo] Confirm Python $\geq$ 3.10: \texttt{python --version}
    \item[\todo] Create virtual environment inside \filepath{python-module/}:
\end{itemize}

\begin{lstlisting}[language=bash]
cd python-module
python -m venv .venv
\end{lstlisting}

\begin{itemize}
    \item[\todo] Create \filepath{pyproject.toml}:
\end{itemize}

\begin{lstlisting}[language=bash]
[project]
name = "helix-python-module"
version = "0.1.0"
description = "Isolation Forest anomaly detection for HELIX"
requires-python = ">=3.10"

[project.dependencies]
scikit-learn = ">=1.3.0"
pandas = ">=2.0.0"
numpy = ">=1.24.0"
joblib = ">=1.3.0"
mysql-connector-python = ">=8.0.0"

[project.optional-dependencies]
dev = [
    "pytest>=7.0.0",
]

[build-system]
requires = ["setuptools>=68.0"]
build-backend = "setuptools.backends.legacy:BuildBackend"
\end{lstlisting}

\begin{itemize}
    \item[\todo] Install: \texttt{pip install -e ".[dev]"} (inside \texttt{.venv})
    \item[\todo] Update \texttt{PYTHON\_PATH} in \filepath{.env} to point to \texttt{.venv} python binary
\end{itemize}

\section{0.11 -- Java Environment}

\begin{itemize}
    \item[\done] Confirm JDK $\geq$ 17: \texttt{java -version}, \texttt{javac -version}
    \item[\todo] Review \filepath{java-module/compile.sh} -- understand what it does:
\end{itemize}

\begin{lstlisting}[language=bash]
#!/bin/bash
mkdir -p build
javac -d build src/LogFormat.java src/ParseResult.java src/LogParser.java src/HashScanner.java
echo "Compilation complete"
\end{lstlisting}

\begin{itemize}
    \item[\todo] Run \filepath{compile.sh} (or equivalent on Windows with \texttt{javac})
    \item[\todo] Verify \texttt{.class} files appear in \filepath{java-module/build/}
\end{itemize}

\section*{Deliverables}

\begin{itemize}
    \item Git repo with clean structure, \filepath{.gitignore}, \filepath{CHANGELOG.md}
    \item Apache vhost serving \filepath{frontend/} at root, \filepath{api/} at \texttt{/api}
    \item Composer installed with PSR-4 autoloading configured
    \item \filepath{api/config/Environment.php}, \filepath{Database.php}, \filepath{Cors.php}
    \item \filepath{api/helpers/Response.php}, \filepath{Validator.php}
    \item \filepath{api/index.php} Front Controller with basic routing skeleton
    \item \filepath{python-module/pyproject.toml} (replacing \filepath{requirements.txt})
    \item Java module compiling without errors
\end{itemize}

\section*{Validation Criteria}

\begin{itemize}
    \item[\todo] \texttt{http://helix.local} $\rightarrow$ serves \filepath{frontend/index.html}
    \item[\todo] \texttt{http://helix.local/api} $\rightarrow$ JSON response (not HTML)
    \item[\todo] \texttt{http://helix.local/api/nonexistent} $\rightarrow$ \texttt{\{"success":false,"error":"Route not found"\}}
    \item[\todo] \texttt{composer dump-autoload} runs without errors
    \item[\todo] Python: \texttt{python -c "import sklearn; print('OK')"} inside \texttt{.venv}
    \item[\todo] Java: \texttt{javac} compiles all source files
    \item[\todo] Git log shows meaningful commits
\end{itemize}

\section*{Common Mistakes to Avoid}

\begin{warnbox}[title=Watch Out]
\begin{itemize}
    \item[$\times$] Skipping the Apache Alias directive -- frontend won't be able to call \texttt{/api/*}
    \item[$\times$] Committing \filepath{.env} -- real credentials must stay out of git
    \item[$\times$] Keeping \filepath{requirements.txt} -- you decided to use \filepath{pyproject.toml}
    \item[$\times$] Not setting up PSR-4 autoloading -- you'll write manual \texttt{require} statements everywhere
    \item[$\times$] Writing route logic directly in \filepath{api/index.php} instead of dispatching to controllers
\end{itemize}
\end{warnbox}

% ══════════════════════════════════════════════════════════════════════════
% PHASE 1
% ══════════════════════════════════════════════════════════════════════════
\chapter{PHASE 1 -- Architecture Design, Database Schema \& UML}

\begin{phasebox}{Phase 1: Architecture \& Design}
\textbf{Why this comes before coding:} You already have \texttt{.puml} diagram files. This phase is about validating them, locking the database schema, and mapping every SRS requirement to a concrete implementation path. Changing the schema after writing 5 models is expensive.
\end{phasebox}

\section*{Objective}
Finalize and validate the complete database schema, UML diagrams, API contract, and traceability matrix. Nothing functional gets built in this phase -- only design artifacts.

\section*{Learning Goals}
\begin{itemize}
    \item Understand relational database normalization (1NF, 2NF, 3NF)
    \item Understand how to derive a schema from use cases and actors
    \item Understand REST API contract-first design
    \item Understand what each UML diagram type communicates and to whom
\end{itemize}

\section*{Key Theoretical Concepts}
\begin{itemize}
    \item \textbf{Entity-Relationship vs Class Diagram:} ER models data; Class models code structure
    \item \textbf{Foreign key CASCADE rules:} \texttt{ON DELETE CASCADE} vs \texttt{ON DELETE SET NULL} -- when to use each
    \item \textbf{ENUM vs lookup table:} ENUMs are fast and sufficient for a PoC; lookup tables scale better
    \item \textbf{REST resource naming:} \texttt{/alerts} not \texttt{/getAlerts}; \texttt{/alerts/\{id\}/status} not \texttt{/updateAlertStatus}
    \item \textbf{HTTP status codes:} 200 OK, 201 Created, 400 Bad Request, 401 Unauthorized, 403 Forbidden, 404 Not Found, 409 Conflict, 422 Unprocessable Entity, 500 Internal Server Error
\end{itemize}

\section{1.1 -- Database Schema}

\filepath{db/schema.sql} -- finalize all tables before writing a single model:

\begin{lstlisting}[language=SQL]
CREATE DATABASE IF NOT EXISTS helix_db CHARACTER SET utf8mb4 COLLATE utf8mb4_unicode_ci;
USE helix_db;

CREATE TABLE users (
    id            INT AUTO_INCREMENT PRIMARY KEY,
    username      VARCHAR(50)  UNIQUE NOT NULL,
    email         VARCHAR(100) UNIQUE NOT NULL,
    password_hash VARCHAR(255) NOT NULL,
    role          ENUM('administrator','blue_team','red_team','purple_team','learner') NOT NULL,
    is_active     BOOLEAN      DEFAULT TRUE,
    created_at    TIMESTAMP    DEFAULT CURRENT_TIMESTAMP,
    last_login    TIMESTAMP    NULL
);

CREATE TABLE log_entries (
    id          INT AUTO_INCREMENT PRIMARY KEY,
    filename    VARCHAR(255),
    source_ip   VARCHAR(45),
    timestamp   TIMESTAMP    NOT NULL,
    severity    ENUM('LOW','MEDIUM','HIGH','CRITICAL') DEFAULT 'LOW',
    event_type  VARCHAR(100),
    message     TEXT         NOT NULL,
    raw_line    TEXT,
    ingested_at TIMESTAMP    DEFAULT CURRENT_TIMESTAMP,
    ingested_by INT          NULL,
    FOREIGN KEY (ingested_by) REFERENCES users(id) ON DELETE SET NULL
);

CREATE TABLE anomaly_scores (
    id            INT AUTO_INCREMENT PRIMARY KEY,
    log_entry_id  INT   NOT NULL,
    score         FLOAT NOT NULL,
    severity_label ENUM('LOW','MEDIUM','HIGH','CRITICAL') NOT NULL,
    computed_at   TIMESTAMP DEFAULT CURRENT_TIMESTAMP,
    FOREIGN KEY (log_entry_id) REFERENCES log_entries(id) ON DELETE CASCADE
);

CREATE TABLE alerts (
    id               INT AUTO_INCREMENT PRIMARY KEY,
    log_entry_id     INT  NULL,
    anomaly_score_id INT  NULL,
    title            VARCHAR(255) NOT NULL,
    description      TEXT,
    severity         ENUM('LOW','MEDIUM','HIGH','CRITICAL') NOT NULL,
    status           ENUM('open','investigating','resolved') DEFAULT 'open',
    assigned_to      INT  NULL,
    created_at       TIMESTAMP DEFAULT CURRENT_TIMESTAMP,
    updated_at       TIMESTAMP DEFAULT CURRENT_TIMESTAMP ON UPDATE CURRENT_TIMESTAMP,
    FOREIGN KEY (log_entry_id)     REFERENCES log_entries(id)    ON DELETE SET NULL,
    FOREIGN KEY (anomaly_score_id) REFERENCES anomaly_scores(id) ON DELETE SET NULL,
    FOREIGN KEY (assigned_to)      REFERENCES users(id)          ON DELETE SET NULL
);

CREATE TABLE chat_sessions (
    id           INT AUTO_INCREMENT PRIMARY KEY,
    user_id      INT  NOT NULL,
    alert_id     INT  NULL,
    message_role ENUM('user','assistant') NOT NULL,
    content      TEXT NOT NULL,
    sent_at      TIMESTAMP DEFAULT CURRENT_TIMESTAMP,
    FOREIGN KEY (user_id)  REFERENCES users(id)  ON DELETE CASCADE,
    FOREIGN KEY (alert_id) REFERENCES alerts(id) ON DELETE SET NULL
);

CREATE TABLE scan_results (
    id           INT AUTO_INCREMENT PRIMARY KEY,
    performed_by INT NOT NULL,
    target       VARCHAR(255) NOT NULL,
    port         INT NOT NULL,
    protocol     VARCHAR(10),
    status       ENUM('open','closed','filtered') NOT NULL,
    scanned_at   TIMESTAMP DEFAULT CURRENT_TIMESTAMP,
    FOREIGN KEY (performed_by) REFERENCES users(id)
);
\end{lstlisting}

\begin{itemize}
    \item[\todo] Verify schema matches your UML diagrams
    \item[\todo] Import: \texttt{mysql -u helix\_user -p helix\_db < db/schema.sql}
    \item[\todo] Confirm all 6 tables appear in phpMyAdmin
\end{itemize}

\textbf{Note on migrations:} Every schema change after the initial import must be a new migration file: \texttt{002\_}, \texttt{003\_}, etc. Never re-edit \filepath{schema.sql} to fix an existing column -- add a migration.

\section{1.2 -- API Route Contract}

Document every route in \filepath{docs/SDD/API\_REFERENCE.md}. For each route define: method, path, required role, request body, success response, error responses.

\textbf{Complete route table:}

\begin{center}
\begin{tabular}{C{1.5cm} L{3.5cm} L{3.5cm} L{3cm} C{1.5cm}}
    \toprule
    \textbf{Method} & \textbf{Path} & \textbf{Controller.Method} & \textbf{Required Role} & \textbf{FR} \\
    \midrule
    POST & /auth/register & AuthController.register & public & FR-007 \\
    POST & /auth/login & AuthController.login & public & FR-001 \\
    POST & /auth/logout & AuthController.logout & any authenticated & FR-005 \\
    GET & /auth/me & AuthController.me & any authenticated & FR-003 \\
    GET & /dashboard/stats & DashboardController.stats & any authenticated & FR-010 \\
    POST & /logs/upload & LogController.upload & admin, blue, purple & FR-020 \\
    GET & /logs & LogController.index & any authenticated & FR-022 \\
    GET & /logs/\{id\} & LogController.show & any authenticated & FR-023 \\
    GET & /alerts & AlertController.index & any authenticated & FR-030 \\
    GET & /alerts/\{id\} & AlertController.show & any authenticated & FR-034 \\
    PUT & /alerts/\{id\}/status & AlertController.updateStatus & admin, blue, purple & FR-033 \\
    POST & /ml/score & MLController.score & administrator & FR-040 \\
    GET & /ml/scores & MLController.index & any authenticated & FR-044 \\
    POST & /ai/chat & AIController.chat & any authenticated & FR-050 \\
    GET & /ai/history & AIController.history & any authenticated & FR-053 \\
    POST & /scanner/scan & ScannerController.scan & admin, red, purple & FR-060 \\
    GET & /scanner/results & ScannerController.results & admin, red, purple & FR-062 \\
    \bottomrule
\end{tabular}
\end{center}

\begin{itemize}
    \item[\todo] Add this table to \filepath{docs/SDD/API\_REFERENCE.md}
    \item[\todo] Add every route to your \filepath{api/index.php} routing table (as empty stubs for now)
\end{itemize}

\section{1.3 -- UML Diagram Validation}

You already have \texttt{.puml} files. This task is about reviewing them for correctness:

\begin{itemize}
    \item[\todo] \filepath{use\_case.puml} -- Does it show all 7 actors (A0--A7)? Are \texttt{<<extend>>} / \texttt{<<include>>} used correctly?
    \item[\todo] \filepath{er\_diagram.puml} -- Does it match \filepath{schema.sql} exactly? Check column names, types, FK arrows
    \item[\todo] \filepath{class\_services.puml} -- Does it show Controller $\rightarrow$ Service $\rightarrow$ Model relationships?
    \item[\todo] \filepath{sequence\_auth.puml} -- Does it show: browser $\rightarrow$ index.php $\rightarrow$ AuthMiddleware $\rightarrow$ AuthController $\rightarrow$ AuthService $\rightarrow$ User model $\rightarrow$ MySQL?
    \item[\todo] \filepath{sequence\_log\_ingest.puml} -- Does it show the pipeline: upload $\rightarrow$ LogService $\rightarrow$ JavaBridge $\rightarrow$ LogParser $\rightarrow$ LogEntry model $\rightarrow$ MLService trigger?
    \item[\todo] \filepath{sequence\_ai\_chat.puml} -- Does it show: AIController $\rightarrow$ AIService $\rightarrow$ OpenAI API $\rightarrow$ ChatSession model?
    \item[\todo] \filepath{component.puml} -- Does it show: Browser, Apache (frontend + api), PHP (controllers/services/models), Python ML, Java module, MySQL?
    \item[\todo] \filepath{activity\_alert.puml} -- Does it show the full alert lifecycle: log ingested $\rightarrow$ scored $\rightarrow$ alert created $\rightarrow$ triaged $\rightarrow$ resolved?
\end{itemize}

Export all diagrams as PNG into \filepath{docs/UML/png/}.

\section{1.4 -- SRS Traceability Matrix}

\filepath{docs/traceability\_matrix.md}:

\begin{center}
\begin{tabular}{L{1.8cm} L{3.5cm} L{3cm} L{3.5cm} L{1.8cm} C{1.2cm}}
    \toprule
    \textbf{FR ID} & \textbf{Description} & \textbf{Route} & \textbf{Controller.Method} & \textbf{DB Table} & \textbf{Status} \\
    \midrule
    FR-001 & Login & POST /auth/login & AuthController.login & users & [\todo] \\
    FR-002 & Failed login handling & POST /auth/login & AuthController.login & users & [\todo] \\
    FR-003 & Session persistence & GET /auth/me & AuthController.me & users & [\todo] \\
    \midrule
    \multicolumn{6}{c}{\textit{... (map every FR-001 to FR-063)}} \\
    \bottomrule
\end{tabular}
\end{center}

\begin{itemize}
    \item[\todo] Map every FR-001 to FR-063 to a route, controller method, and DB table
    \item[\todo] This matrix is your implementation checklist -- check FRs off as you build
\end{itemize}

\section{1.5 -- Seed File}

\filepath{database/seed.sql}:

\begin{lstlisting}[language=SQL]
-- Admin user (password: Admin@1234)
-- Generate hash: php -r "echo password_hash('Admin@1234', PASSWORD_BCRYPT, ['cost'=>12]);"
INSERT INTO users (username, email, password_hash, role) VALUES
('admin', 'admin@helix.local', '$2y$12$GENERATED_HASH_HERE', 'administrator'),
('blue1', 'blue@helix.local',  '$2y$12$GENERATED_HASH_HERE', 'blue_team'),
('red1',  'red@helix.local',   '$2y$12$GENERATED_HASH_HERE', 'red_team');
\end{lstlisting}

\begin{itemize}
    \item[\todo] Generate real bcrypt hashes and update seed.sql
    \item[\todo] Import: \texttt{mysql -u helix\_user -p helix\_db < database/seed.sql}
\end{itemize}

\section*{Deliverables}

\begin{itemize}
    \item \filepath{db/schema.sql} -- final, importable, 6 tables
    \item \filepath{docs/SDD/API\_REFERENCE.md} -- complete route contract
    \item All UML diagrams validated and exported to PNG
    \item \filepath{docs/traceability\_matrix.md} -- all FRs mapped
    \item \filepath{database/seed.sql} -- 3 test users
\end{itemize}

\section*{Validation Criteria}

\begin{itemize}
    \item[\todo] Schema imports without error
    \item[\todo] All 6 tables visible in phpMyAdmin
    \item[\todo] All UML diagrams reviewed and corrected
    \item[\todo] All FR-001 to FR-063 appear in traceability matrix
    \item[\todo] All routes registered as stubs in \filepath{api/index.php}
    \item[\todo] Git commit: \texttt{docs: finalize schema, API contract, UML and traceability matrix}
\end{itemize}

\section*{Common Mistakes to Avoid}

\begin{warnbox}[title=Watch Out]
\begin{itemize}
    \item[$\times$] Editing \filepath{schema.sql} after first import instead of creating a migration
    \item[$\times$] Skipping the traceability matrix -- you will lose track of which FRs are done
    \item[$\times$] Routes in \filepath{api/index.php} pointing to non-existent controller methods
    \item[$\times$] UML diagrams that don't match the actual code structure
\end{itemize}
\end{warnbox}

% ══════════════════════════════════════════════════════════════════════════
% PHASE 2
% ══════════════════════════════════════════════════════════════════════════
\chapter{PHASE 2 -- Authentication Module (FR-001 to FR-007)}

\begin{phasebox}{Phase 2: Authentication \& RBAC}
\textbf{Why this comes before all other modules:} Every other feature is gated behind authentication. You cannot build role-based features without a working session system. Auth is also your highest-security concern -- if it's broken, everything built on it is insecure.
\end{phasebox}

\section*{Objective}
Implement the complete authentication flow: register, login, logout, session validation -- with bcrypt hashing, PHP sessions, and role middleware enforcing access control on every protected route.

\section*{Learning Goals}
\begin{itemize}
    \item Understand bcrypt and why password hashing $\neq$ encryption
    \item Understand PHP session security: \texttt{session\_regenerate\_id()}, cookie flags
    \item Understand the difference between Authentication (who are you?) and Authorization (what can you do?)
    \item Understand how middleware chains work in a Front Controller pattern
\end{itemize}

\section*{Key Theoretical Concepts}
\begin{itemize}
    \item \textbf{bcrypt cost factor:} Each increment doubles computation time -- cost 12 is $\sim$300ms, which is too slow to brute-force
    \item \textbf{Session fixation attack:} Why \texttt{session\_regenerate\_id(true)} is mandatory after login
    \item \textbf{HttpOnly cookie flag:} Prevents JavaScript from reading the session cookie (XSS protection)
    \item \textbf{SameSite=Strict:} Prevents the browser from sending the cookie on cross-site requests (CSRF protection)
    \item \textbf{Generic error messages:} "Invalid credentials" -- never "User not found" or "Wrong password"
\end{itemize}

\section{2.1 -- Auth Middleware}

\filepath{api/middleware/AuthMiddleware.php}:

\begin{lstlisting}[language=PHP]
<?php
namespace App\Middleware;

use App\Helpers\Response;

class AuthMiddleware
{
    public static function handle(): void
    {
        self::startSecureSession();

        if (!isset($_SESSION['user_id'])) {
            Response::unauthorized();
        }
    }

    public static function startSecureSession(): void
    {
        if (session_status() === PHP_SESSION_NONE) {
            ini_set('session.cookie_httponly', '1');
            ini_set('session.cookie_samesite', 'Strict');
            ini_set('session.gc_maxlifetime', (string)($_ENV['SESSION_LIFETIME'] ?? 3600));
            session_start();
        }
    }

    public static function currentUser(): array
    {
        return [
            'id'   => $_SESSION['user_id']   ?? null,
            'role' => $_SESSION['user_role'] ?? null,
        ];
    }
}
\end{lstlisting}

\filepath{api/middleware/RoleMiddleware.php}:

\begin{lstlisting}[language=PHP]
<?php
namespace App\Middleware;

use App\Helpers\Response;

class RoleMiddleware
{
    public static function require(array $allowedRoles): void
    {
        AuthMiddleware::handle();

        $userRole = $_SESSION['user_role'] ?? '';

        if (!in_array($userRole, $allowedRoles, true)) {
            Response::forbidden();
        }
    }
}
\end{lstlisting}

\section{2.2 -- User Model}

\filepath{api/models/User.php}:

\begin{lstlisting}[language=PHP]
<?php
namespace App\Models;

use App\Config\Database;
use PDO;

class User
{
    public static function findByUsername(string $username): ?array
    {
        $db   = Database::getInstance();
        $stmt = $db->prepare('SELECT * FROM users WHERE username = ? AND is_active = 1 LIMIT 1');
        $stmt->execute([$username]);
        return $stmt->fetch() ?: null;
    }

    public static function findById(int $id): ?array
    {
        $db   = Database::getInstance();
        $stmt = $db->prepare('SELECT id, username, email, role, created_at, last_login FROM users WHERE id = ? LIMIT 1');
        $stmt->execute([$id]);
        return $stmt->fetch() ?: null;
    }

    public static function create(string $username, string $email, string $passwordHash, string $role = 'learner'): int
    {
        $db   = Database::getInstance();
        $stmt = $db->prepare('INSERT INTO users (username, email, password_hash, role) VALUES (?, ?, ?, ?)');
        $stmt->execute([$username, $email, $passwordHash, $role]);
        return (int) $db->lastInsertId();
    }

    public static function updateLastLogin(int $id): void
    {
        $db   = Database::getInstance();
        $stmt = $db->prepare('UPDATE users SET last_login = NOW() WHERE id = ?');
        $stmt->execute([$id]);
    }

    public static function emailOrUsernameExists(string $username, string $email): bool
    {
        $db   = Database::getInstance();
        $stmt = $db->prepare('SELECT id FROM users WHERE username = ? OR email = ? LIMIT 1');
        $stmt->execute([$username, $email]);
        return (bool) $stmt->fetch();
    }
}
\end{lstlisting}

\section{2.3 -- Auth Service}

\filepath{api/services/AuthService.php}:

\begin{lstlisting}[language=PHP]
<?php
namespace App\Services;

use App\Models\User;
use App\Middleware\AuthMiddleware;

class AuthService
{
    public function register(array $data): array
    {
        if (User::emailOrUsernameExists($data['username'], $data['email'])) {
            return ['error' => 'Username or email already exists', 'code' => 409];
        }

        $hash = password_hash($data['password'], PASSWORD_BCRYPT, ['cost' => 12]);
        $id   = User::create($data['username'], $data['email'], $hash);

        return ['user_id' => $id];
    }

    public function login(string $username, string $password): array
    {
        $user = User::findByUsername($username);

        if (!$user || !password_verify($password, $user['password_hash'])) {
            return ['error' => 'Invalid credentials', 'code' => 401];
        }

        AuthMiddleware::startSecureSession();
        session_regenerate_id(true);

        $_SESSION['user_id']   = $user['id'];
        $_SESSION['user_role'] = $user['role'];

        User::updateLastLogin($user['id']);

        return [
            'user' => [
                'id'       => $user['id'],
                'username' => $user['username'],
                'role'     => $user['role'],
            ]
        ];
    }

    public function logout(): void
    {
        AuthMiddleware::startSecureSession();
        $_SESSION = [];
        session_destroy();
    }

    public function currentUser(): ?array
    {
        $id = $_SESSION['user_id'] ?? null;
        return $id ? User::findById((int)$id) : null;
    }
}
\end{lstlisting}

\section{2.4 -- Auth Controller}

\filepath{api/controllers/AuthController.php}:

\begin{lstlisting}[language=PHP]
<?php
namespace App\Controllers;

use App\Services\AuthService;
use App\Middleware\AuthMiddleware;
use App\Helpers\Response;
use App\Helpers\Validator;

class AuthController
{
    private AuthService $service;

    public function __construct()
    {
        $this->service = new AuthService();
    }

    public function register(): void
    {
        $body = json_decode(file_get_contents('php://input'), true) ?? [];

        $v = (new Validator())
            ->required('username', $body['username'] ?? '')
            ->required('email', $body['email'] ?? '')
            ->required('password', $body['password'] ?? '')
            ->email('email', $body['email'] ?? '')
            ->minLength('password', $body['password'] ?? '', 8);

        if (!$v->passes()) {
            Response::error(json_encode($v->errors()), 422);
        }

        $result = $this->service->register($body);

        if (isset($result['error'])) {
            Response::error($result['error'], $result['code']);
        }

        Response::success($result, 'Account created', 201);
    }

    public function login(): void
    {
        $body = json_decode(file_get_contents('php://input'), true) ?? [];

        $result = $this->service->login(
            $body['username'] ?? '',
            $body['password'] ?? ''
        );

        if (isset($result['error'])) {
            Response::error($result['error'], $result['code']);
        }

        Response::success($result['user'], 'Login successful');
    }

    public function logout(): void
    {
        AuthMiddleware::handle();
        $this->service->logout();
        Response::success(null, 'Logged out');
    }

    public function me(): void
    {
        AuthMiddleware::handle();
        $user = $this->service->currentUser();

        if (!$user) {
            Response::unauthorized();
        }

        Response::success($user);
    }
}
\end{lstlisting}

\section{2.5 -- Register Auth Routes}

In \filepath{api/index.php}, add to the routes array:

\begin{lstlisting}[language=PHP]
use App\Controllers\AuthController;

$routes['POST']['/auth/register'] = [AuthController::class, 'register'];
$routes['POST']['/auth/login']    = [AuthController::class, 'login'];
$routes['POST']['/auth/logout']   = [AuthController::class, 'logout'];
$routes['GET']['/auth/me']        = [AuthController::class, 'me'];
\end{lstlisting}

\section{2.6 -- Frontend: Login Page}

\filepath{frontend/login.html} -- already exists in your tree. Build it now:
\begin{itemize}
    \item HTML form: username + password fields
    \item No page refresh -- use JavaScript \texttt{fetch()}
\end{itemize}

\filepath{frontend/js/auth.js}:

\begin{lstlisting}[language=JavaScript]
async function login(username, password) {
  const res = await fetch("/api/auth/login", {
    method: "POST",
    headers: { "Content-Type": "application/json" },
    credentials: "include",
    body: JSON.stringify({ username, password }),
  });

  const data = await res.json();

  if (data.success) {
    sessionStorage.setItem("userRole", data.data.role);
    window.location.href = "/dashboard.html";
  } else {
    showError(data.error);
  }
}
\end{lstlisting}

\filepath{frontend/js/api.js} -- the shared API client module:

\begin{lstlisting}[language=JavaScript]
const API_BASE = "/api";

async function apiFetch(path, options = {}) {
  const defaults = {
    credentials: "include",
    headers: { "Content-Type": "application/json" },
  };

  const res = await fetch(`${API_BASE}${path}`, { ...defaults, ...options });

  if (res.status === 401) {
    window.location.href = "/login.html";
    return;
  }

  return await res.json();
}

const api = {
  get: (path) => apiFetch(path),
  post: (path, body) =>
    apiFetch(path, { method: "POST", body: JSON.stringify(body) }),
  put: (path, body) =>
    apiFetch(path, { method: "PUT", body: JSON.stringify(body) }),
  delete: (path) => apiFetch(path, { method: "DELETE" }),
};
\end{lstlisting}

\textbf{Why \texttt{credentials: 'include'}:} Without this, the browser won't send the session cookie with cross-path requests. This is the most common cause of "keeps logging me out" bugs.

\section{2.7 -- Auth Tests}

\filepath{tests/php/AuthServiceTest.php} (PHPUnit):

\begin{lstlisting}[language=PHP]
<?php
namespace Tests;

use PHPUnit\Framework\TestCase;
use App\Services\AuthService;

class AuthServiceTest extends TestCase
{
    public function test_register_returns_user_id(): void
    {
        $this->markTestIncomplete('Requires test DB setup');
    }

    public function test_login_returns_error_on_wrong_password(): void
    {
        $service = new AuthService();
        $result  = $this->invokeMethod($service, 'validateCredentials', ['nonexistent', 'wrongpass']);
        $this->assertArrayHasKey('error', $result);
    }

    private function invokeMethod(object $obj, string $method, array $args = []): mixed
    {
        $ref = new \ReflectionMethod($obj, $method);
        $ref->setAccessible(true);
        return $ref->invokeArgs($obj, $args);
    }
}
\end{lstlisting}

\section{2.8 -- Manual Test Checklist}

\filepath{tests/manual/auth\_tests.md}:

\begin{itemize}
    \item[\todo] POST \texttt{/api/auth/register} valid data $\rightarrow$ 201, user\_id returned
    \item[\todo] POST \texttt{/api/auth/register} duplicate username $\rightarrow$ 409
    \item[\todo] POST \texttt{/api/auth/register} invalid email $\rightarrow$ 422
    \item[\todo] POST \texttt{/api/auth/login} valid credentials $\rightarrow$ 200, role returned, cookie set
    \item[\todo] POST \texttt{/api/auth/login} wrong password $\rightarrow$ 401, "Invalid credentials"
    \item[\todo] GET \texttt{/api/auth/me} with cookie $\rightarrow$ 200, user data
    \item[\todo] GET \texttt{/api/auth/me} without cookie $\rightarrow$ 401
    \item[\todo] POST \texttt{/api/auth/logout} $\rightarrow$ 200, cookie destroyed
    \item[\todo] GET \texttt{/api/auth/me} after logout $\rightarrow$ 401
    \item[\todo] Test route requiring admin role as learner $\rightarrow$ 403
\end{itemize}

Test all using Postman or Insomnia -- configure "Send cookies" option.

\section*{Deliverables}

\begin{itemize}
    \item \filepath{api/middleware/AuthMiddleware.php}, \filepath{RoleMiddleware.php}
    \item \filepath{api/models/User.php}
    \item \filepath{api/services/AuthService.php}
    \item \filepath{api/controllers/AuthController.php}
    \item \filepath{frontend/login.html} functional form
    \item \filepath{frontend/js/auth.js}, \filepath{frontend/js/api.js}
    \item Auth routes registered in \filepath{api/index.php}
    \item \filepath{tests/manual/auth\_tests.md} all checks passing
\end{itemize}

\section*{Validation Criteria}

\begin{itemize}
    \item[\todo] FR-001 to FR-007 checked in traceability matrix
    \item[\todo] Passwords stored as bcrypt (verify in phpMyAdmin -- hash starts with \texttt{\$2y\$})
    \item[\todo] Wrong password returns same error message as non-existent user
    \item[\todo] \texttt{session\_regenerate\_id(true)} called on every login
    \item[\todo] Frontend redirects to login on 401 (handled in \filepath{api.js})
    \item[\todo] Git commit: \texttt{feat(auth): implement authentication, RBAC middleware, login frontend}
\end{itemize}

\section*{Common Mistakes to Avoid}

\begin{warnbox}[title=Watch Out]
\begin{itemize}
    \item[$\times$] Checking roles only in the frontend JavaScript -- always enforce in \texttt{RoleMiddleware}
    \item[$\times$] Returning "User not found" vs "Wrong password" separately -- merge into one message
    \item[$\times$] Forgetting \texttt{credentials: 'include'} in every \texttt{fetch()} call
    \item[$\times$] Using \texttt{md5()} or \texttt{sha1()} for passwords
    \item[$\times$] Not calling \texttt{session\_regenerate\_id()} -- session fixation vulnerability
\end{itemize}
\end{warnbox}

% ══════════════════════════════════════════════════════════════════════════
% PHASE 3
% ══════════════════════════════════════════════════════════════════════════
\chapter{PHASE 3 -- Dashboard Module (FR-010 to FR-014)}

\begin{phasebox}{Phase 3: Dashboard \& Routing}
\textbf{Why now:} The dashboard is the shell of the application -- the first page after login. Building it now, even with placeholder data, establishes the UI framework, client-side router, and nav structure that all subsequent modules will plug into.
\end{phasebox}

\section*{Objective}
Build the dashboard frontend shell (with routing), implement the stats API endpoint, and render security metrics with Chart.js. Use placeholder/zero data for sections not yet implemented.

\section*{Learning Goals}
\begin{itemize}
    \item Understand client-side routing without a framework (\filepath{router.js})
    \item Understand the HTML5 History API (\texttt{pushState}, \texttt{popstate})
    \item Understand Chart.js configuration and data binding
    \item Understand role-based UI rendering in JavaScript
\end{itemize}

\section*{Key Theoretical Concepts}
\begin{itemize}
    \item \textbf{SPA routing:} The browser loads one HTML page; JavaScript swaps content based on URL without full page reload
    \item \textbf{\texttt{history.pushState()}:} Changes the URL without navigating -- the foundation of client-side routing
    \item \textbf{Dashboard as aggregation endpoint:} The dashboard API doesn't own data; it queries multiple tables and returns a summary
    \item \textbf{Graceful empty states:} Your dashboard must not crash when the DB is empty (it will be during early development)
\end{itemize}

\section{3.1 -- Client-Side Router}

\filepath{frontend/js/router.js}:

\begin{lstlisting}[language=JavaScript]
const routes = {
  "/dashboard": () => import("./dashboard.js").then((m) => m.render()),
  "/logs": () => import("./logs.js").then((m) => m.render()),
  "/alerts": () => import("./alerts.js").then((m) => m.render()),
  "/ai": () => import("./chat.js").then((m) => m.render()),
  "/scanner": () => import("./scanner.js").then((m) => m.render()),
};

async function navigate(path) {
  const handler = routes[path];

  if (!handler) {
    document.getElementById("app").innerHTML = "<h2>404 - Page not found</h2>";
    return;
  }

  history.pushState({}, "", path);
  document.getElementById("app").innerHTML =
    '<div class="loading">Loading...</div>';
  await handler();
}

window.addEventListener("popstate", () => navigate(location.pathname));

document.addEventListener("click", (e) => {
  const link = e.target.closest("[data-route]");
  if (link) {
    e.preventDefault();
    navigate(link.dataset.route);
  }
});

export { navigate };
\end{lstlisting}

\section{3.2 -- Dashboard Shell}

\filepath{frontend/dashboard.html}:

\begin{lstlisting}[language=HTML]
<!DOCTYPE html>
<html lang="en">
  <head>
    <meta charset="UTF-8" />
    <title>HELIX - Dashboard</title>
    <link rel="stylesheet" href="/css/main.css" />
    <link rel="stylesheet" href="/css/dashboard.css" />
  </head>
  <body>
    <div id="layout">
      <nav id="sidebar"></nav>
      <main id="app"></main>
    </div>
    <script type="module" src="/js/ui.js"></script>
    <script type="module" src="/js/router.js"></script>
  </body>
</html>
\end{lstlisting}

\filepath{frontend/js/ui.js} -- builds the sidebar based on user role:

\begin{lstlisting}[language=JavaScript]
import { navigate } from "./router.js";

const roleNavMap = {
  administrator: ["/dashboard", "/logs", "/alerts", "/ai", "/scanner"],
  blue_team: ["/dashboard", "/logs", "/alerts", "/ai"],
  red_team: ["/dashboard", "/scanner", "/ai"],
  purple_team: ["/dashboard", "/logs", "/alerts", "/ai", "/scanner"],
  learner: ["/dashboard", "/ai"],
};

async function initUI() {
  const res = await fetch("/api/auth/me", { credentials: "include" });
  if (!res.ok) {
    window.location.href = "/login.html";
    return;
  }

  const { data: user } = await res.json();
  const allowedPaths = roleNavMap[user.role] ?? [];

  buildSidebar(allowedPaths, user);
  navigate("/dashboard");
}

function buildSidebar(paths, user) {
  const labels = {
    "/dashboard": "Dashboard",
    "/logs": "Logs",
    "/alerts": "Alerts",
    "/ai": "AI Assistant",
    "/scanner": "Scanner",
  };
  const nav = document.getElementById("sidebar");
  nav.innerHTML =
    paths
      .map((p) => `<a data-route="${p}" href="${p}">${labels[p]}</a>`)
      .join("") + `<button id="logout-btn">Logout (${user.username})</button>`;

  document.getElementById("logout-btn").addEventListener("click", async () => {
    await fetch("/api/auth/logout", { method: "POST", credentials: "include" });
    window.location.href = "/login.html";
  });
}

initUI();
\end{lstlisting}

\section{3.3 -- Dashboard Stats Endpoint}

\filepath{api/controllers/DashboardController.php}:

\begin{lstlisting}[language=PHP]
<?php
namespace App\Controllers;

use App\Config\Database;
use App\Middleware\AuthMiddleware;
use App\Helpers\Response;

class DashboardController
{
    public function stats(): void
    {
        AuthMiddleware::handle();
        $db = Database::getInstance();

        $alertsBySeverity = $db->query(
            "SELECT severity, COUNT(*) as count FROM alerts GROUP BY severity"
        )->fetchAll();

        $logCount = $db->query("SELECT COUNT(*) as count FROM log_entries")->fetch()['count'];

        $avgScore = $db->query("SELECT ROUND(AVG(score), 4) as avg FROM anomaly_scores")->fetch()['avg'] ?? 0;

        $trend = $db->query(
            "SELECT DATE(created_at) as date, COUNT(*) as count
             FROM alerts
             WHERE created_at >= DATE_SUB(NOW(), INTERVAL 7 DAY)
             GROUP BY DATE(created_at)
             ORDER BY date ASC"
        )->fetchAll();

        $openAlerts = $db->query(
            "SELECT COUNT(*) as count FROM alerts WHERE status = 'open'"
        )->fetch()['count'];

        Response::success([
            'alerts_by_severity' => $alertsBySeverity,
            'log_count'          => (int) $logCount,
            'avg_anomaly_score'  => (float) $avgScore,
            'alert_trend'        => $trend,
            'open_alerts'        => (int) $openAlerts,
        ]);
    }
}
\end{lstlisting}

Register route in \filepath{api/index.php}:

\begin{lstlisting}[language=PHP]
$routes['GET']['/dashboard/stats'] = [DashboardController::class, 'stats'];
\end{lstlisting}

\section{3.4 -- Dashboard Frontend Module}

\filepath{frontend/js/dashboard.js}:

\begin{lstlisting}[language=JavaScript]
import { api } from "./api.js";

export async function render() {
  const { data } = await api.get("/dashboard/stats");

  document.getElementById("app").innerHTML = `
    <div class="dashboard-grid">
      <div class="metric-card">
        <span class="metric-label">Open Alerts</span>
        <span class="metric-value">${data.open_alerts}</span>
      </div>
      <div class="metric-card">
        <span class="metric-label">Total Logs</span>
        <span class="metric-value">${data.log_count}</span>
      </div>
      <div class="metric-card">
        <span class="metric-label">Avg Anomaly Score</span>
        <span class="metric-value">${data.avg_anomaly_score}</span>
      </div>
      <canvas id="severity-chart"></canvas>
      <canvas id="trend-chart"></canvas>
    </div>
  `;

  renderSeverityChart(data.alerts_by_severity);
  renderTrendChart(data.alert_trend);
}

function renderSeverityChart(alerts) {
  // Chart.js doughnut chart
  const ctx = document.getElementById("severity-chart").getContext("2d");
  new Chart(ctx, {
    type: "doughnut",
    data: {
      labels: alerts.map((a) => a.severity),
      datasets: [{
        data: alerts.map((a) => a.count),
        backgroundColor: ["#dc2626", "#f59e0b", "#3b82f6", "#6b7280"],
      }],
    },
    options: { plugins: { title: { display: true, text: "Alerts by Severity" } } },
  });
}

function renderTrendChart(trend) {
  const ctx = document.getElementById("trend-chart").getContext("2d");
  new Chart(ctx, {
    type: "line",
    data: {
      labels: trend.map((t) => t.date),
      datasets: [{
        label: "Alerts per day",
        data: trend.map((t) => t.count),
        borderColor: "#3C3489",
        tension: 0.3,
      }],
    },
    options: { plugins: { title: { display: true, text: "7-Day Alert Trend" } } },
  });
}
\end{lstlisting}

\section*{Deliverables}

\begin{itemize}
    \item \filepath{frontend/js/router.js}, \filepath{frontend/js/ui.js}
    \item \filepath{frontend/dashboard.html} with sidebar and \texttt{\#app} container
    \item \filepath{frontend/js/dashboard.js} with Chart.js rendering
    \item \filepath{api/controllers/DashboardController.php}
    \item Dashboard route registered in \filepath{api/index.php}
    \item Chart.js included via CDN in \filepath{dashboard.html}
\end{itemize}

\section*{Validation Criteria}

\begin{itemize}
    \item[\todo] FR-010 to FR-014 checked in traceability matrix
    \item[\todo] Sidebar shows only role-appropriate links
    \item[\todo] Dashboard renders with zero data (no crashes)
    \item[\todo] Chart.js charts render with empty/zero datasets
    \item[\todo] Client-side routing works (back/forward buttons)
    \item[\todo] Git commit: \texttt{feat(dashboard): implement dashboard shell, stats API, Chart.js}
\end{itemize}

\section*{Common Mistakes to Avoid}

\begin{warnbox}[title=Watch Out]
\begin{itemize}
    \item[$\times$] Hardcoding role-based nav in HTML -- use \filepath{ui.js} to build dynamically
    \item[$\times$] Dashboard crashing on empty DB -- always use \texttt{?? 0} or \texttt{|| []}
    \item[$\times$] Forgetting \texttt{history.pushState()} -- browser back button breaks
    \item[$\times$] Mixing frontend and backend routing -- keep them separate
\end{itemize}
\end{warnbox}

% ══════════════════════════════════════════════════════════════════════════
% PHASE 4-10 PLACEHOLDER
% ══════════════════════════════════════════════════════════════════════════
\chapter{PHASE 4 through PHASE 10 -- Implementation Phases}

\begin{infobox}[title=Roadmap Continuation]
Phases 4 through 10 follow the same structural pattern as Phases 0--3:
\begin{itemize}
    \item \textbf{Phase 4:} Log Ingestion Module (FR-020 to FR-023)
    \item \textbf{Phase 5:} Alert Module (FR-030 to FR-034)
    \item \textbf{Phase 6:} ML Detection Module (FR-040 to FR-044)
    \item \textbf{Phase 7:} AI Assistant Module (FR-050 to FR-053)
    \item \textbf{Phase 8:} Scanner Module (FR-060 to FR-062)
    \item \textbf{Phase 9:} Integration \& Testing
    \item \textbf{Phase 10:} Docker \& Deployment
\end{itemize}
Each phase includes: Objective, Learning Goals, Key Theoretical Concepts, Practical Tasks with code, Deliverables, Validation Criteria, and Common Mistakes to Avoid.
\end{infobox}

% ══════════════════════════════════════════════════════════════════════════
% CONCLUSION
% ══════════════════════════════════════════════════════════════════════════
\chapter*{Conclusion}
\addcontentsline{toc}{chapter}{Conclusion}

\begin{infobox}[title=Summary]
This roadmap provides a phased, dependency-ordered execution plan for building the HELIX PoC. Each phase is self-contained with clear deliverables, validation criteria, and common pitfalls to avoid. Follow the strict dependency chain -- do not skip ahead. Build one module at a time, test it thoroughly, then proceed.
\end{infobox}

\end{document}
